\section{Conclusion}
\label{sec:conclusion}

\dnote{Three versions of this section are drafted below, one per outcome band
of \code{../OUTCOMES.md}. Keep exactly one before submission and delete the
other two together with this note. Drafting them in advance is deliberate: the
point of the pre-registration is that the wording does not get chosen after the
numbers arrive.}

% ===========================================================================
% VERSION A -- constructive (green / amber-green band)
% ===========================================================================

Allocating a spherical-harmonic degree budget from the local force error
optimizes the wrong functional. Written in the form the dynamics impose, the
objective takes the norm after a signed integral against the state-transition
matrix, which couples every pair of epochs and makes the allocation
non-separable. The classical multiplier treatment and its sensitivity-weighted
repair both still \emph{work} --- they are separable and a single multiplier
decouples them --- but they optimize the wrong functional, and the discretized
form says exactly how: all three criteria are readings of one trajectory
kernel, which the force-level rule ignores, the sensitivity-weighted repair
reduces to its local diagonal, and only the third keeps whole. The weight was
never the missing ingredient, because the weight \emph{is} that diagonal.

The coupling nonetheless has a prefix--suffix structure that keeps a full sweep
over the decision intervals at $\mathcal{O}(M\lvert\mathcal{N}\rvert)$ and
admits a Frank--Wolfe bound on its convex-hull relaxation, so the problem is
solved and certified without propagating anything, at a certified error gap of
\ph{pct}. On \ph{n} orbits of \ph{n} independent designs the certified
allocation is a median \ph{value} better than the equal-budget constant degree
and \ph{value} better than the strongest previously available deployable
schedule, and \ph{pct} of that advantage lives off the diagonal.

Carrying it into a deployable rule turns on a quantity a tail criterion does
not provide: the \emph{direction} of the omitted acceleration. That direction
varies on the field's own resolution scale $\pi r/N$, which is \ph{value}~km at
degree 300, so a surrogate coarser than that scale cannot represent it and one
fine enough to do so carries a number of degrees of freedom comparable to the
omitted coefficients themselves. It can be probed, and the probe is the
controller's dominant cost rather than a rounding term: covering a decision
interval takes a measured \ph{pct} of the total budget, and recovers
the direction to $\kappa=\ph{value}$. The same length scale sets how far ahead
the probe may look and how long a plan made on a cheap pilot arc stays valid,
and those two answers differ: forward over one decision interval, yes; across a
seven-day arc, no. That is why the controller plans its cancellation target
offline and probes forward one interval at a time rather than uploading a
schedule.

Under a common realized-work ceiling, with its own overhead charged in full,
that controller has the smaller seven-day position error than the constant degree on
\ph{n} of \ph{n} resolved comparisons across \ph{n} populations, capturing a
median \ph{value} of the certified benchmark's advantage. \ph{state the
$\beta=0.5$ outcome explicitly.} \ph{state the wide-elliptic outcome
explicitly.}

Every conclusion is model-relative and bounded by the tested designs, arc
lengths and budgets. \ph{closing sentence naming the sharpest bound.}

% ===========================================================================
% VERSION B -- benchmark real, controller cannot capture it (amber band)
% ===========================================================================
%
% \ph{Draft:} the formulation and the certificate stand; the benchmark shows a
% median \ph{value} of attainable gain over the constant degree; no deployable
% controller captures a resolvable fraction of it. The decorrelation argument
% then carries the paper: what separates the benchmark from every deployable
% rule is not the weighting and not the algorithm but access to the omitted
% direction at the field's own resolution scale, and the probe measures exactly
% how much of that access is affordable. State the capture fraction, name which
% substitution destroys it (tail amplitude, pilot arc, low-degree $\Phi$, or
% probe depth), and present the result as an information-limit statement that
% strengthens rather than contradicts the previous paper.
%
% ===========================================================================
% VERSION C -- no attainable gain (orange band)
% ===========================================================================
%
% \ph{Draft:} even with the sensitivity and the sign, and with a certificate
% bounding how much better any schedule of this class could be, allocation does
% not systematically improve on a constant degree at equal realized work in
% this regime. That closes the question the previous paper left open rather
% than leaving it open, and it is worth stating without hedging. Distinguish
% the two ways it can happen --- no gain in the linearized problem at all, or a
% linearized gain that does not survive propagation
% (Section~\ref{sec:campaign-fixedpoint}) --- because they mean different
% things: the first says the objective has no exploitable structure at these
% budgets, the second says the structure is an artefact of optimizing along a
% fixed reference arc.
